\documentclass[12pt]{article}
\usepackage[utf8]{inputenc}
\usepackage[T1]{fontenc}
\usepackage{hyperref}
\title{Xakriabá Language Resources for K-12 Curriculum Development}
\author{Compiled by Claude Sonnet 5.5 (claude-sonnet-5-5) for Andr\'as Kornai}
\date{}
\begin{document}
\maketitle

\section{Introduction}
Xakriabá (ISO 639-3 xkr; Glottocode xakr1238) is a dormant Central Jê language formerly spoken in Minas Gerais, Brazil. The Xakriabá people now speak Portuguese. The UNESCO Atlas lists the language as extinct, and Wikipedia says it is known only from two short wordlists collected by Saint-Hilaire and Eschwege. The community uses the name Akwẽ for its ancestral language and is working to reclaim it, mainly through songs.

\section{Primary Educational Resources}
No primers or textbooks for Xakriabá were found. A 2026 news report describes a digital TV project, Brasil Antenado, that has begun translating informational materials into Akwẽ as a revitalization tool. Whether this counts as a school programme was not established.

\section{Linguistic and Reference Materials}
\begin{itemize}
\item Wikipedia, ``Xakriabá language''. \url{https://en.wikipedia.org/wiki/Xakriab%C3%A1_language}
\item Rodrigues and Silva (2024), essay in \emph{Cadernos de Linguística} on strategies for reclaiming Akwẽ Xakriabá through songs. \url{https://cadernos.abralin.org/index.php/cadernos/article/download/691/900/13606}
\item ``Gest\~ao territorial na Terra Ind\'igena Xakriab\'a e a geopol\'itica das retomadas'' (territorial management and community context). \url{https://www.revistas.una.ac.cr/index.php/geografica/article/view/3022}
\end{itemize}

\section{Religious and Community Resources}
No religious resources were found. Community context appears in the articles above.

\section{Digital and Audio Resources}
\begin{itemize}
\item Estado de Minas (2026), report on technology as a tool against the erasure of the Xakriabá language. \url{https://www.em.com.br/gerais/2026/04/amp/7401314-mg-tecnologia-vira-ferramenta-contra-o-apagamento-da-lingua-dos-xakriabas.html}
\item Correio Braziliense (2026), report on the same topic. \url{https://www.correiobraziliense.com.br/brasil/2026/04/7401337-tecnologia-vira-ferramenta-contra-o-apagamento-da-lingua-dos-xakriabas.html}
\end{itemize}
A site on the International Phonetic Alphabet has a Xakriabá page, but it is generic and gives no verifiable detail, so it is not recommended.

\section{Curriculum Development Context}
This is a heritage and revitalization setting with almost no recorded linguistic material. Song-based work led by the community is the most promising route, and community leadership should guide any curriculum. Because the documented base is so small, the curriculum should be honest that reconstruction is limited, and it should rely on community-held knowledge and on Portuguese-medium indigenous education frameworks.

\end{document}